% Nguồn: decision log 2026-11-15 16:30 (Dev v4, Test-4), 17:30 (24 ca thích nghi), 18:20 (niêm phong, lực, H1--H4),
% 2026-11-17 19:25 (final_config_v5), 2026-11-18 19:30--20:55 (Gemini, thích nghi v3, AgentDojo v3), 2026-11-19 00:25
% (baseline bổ sung, khai báo trước, khám phá). Số: facts.md mục "Bổ sung cho bài v3" và mục Test-4.
% Nguồn gốc Test-4: sinh từ spec, owner và Minh rà soát (không ghi "do owner và Minh soạn", không báo kappa).
% 20/11 (soát văn phong, giữ 15 trang): bỏ câu kiểm hash mỗi lần chạy (đã có ý niêm phong).
% 20/11 (lượt 3, theo bài tham khảo): câu mở mục, bớt ngoặc có ký hiệu, Gemini chỉ ghi là mô hình đóng của Google.
\section{Dữ liệu và giao thức đánh giá}\label{sec:protocol}

Phần đánh giá kiểm tra \ours{} trên dữ liệu chưa dùng khi thiết kế, so với các phòng thủ đã công bố và với các cấu hình
luật đơn giản hơn tại broker.

\paragraph{Tập dữ liệu.} Bộ dữ liệu VN-AgentInject hiện thực hóa mô hình đe dọa ở Mục~\ref{sec:threat} thành các ca chạy
được với dữ liệu tổng hợp. \Dev{} gồm 168 họ và 1.960 ca (840 tấn công, 840 \bPlain, 280 \bAuth), hợp từ tập phát triển
ban đầu 48 họ và hai tập kiểm thử 72 và 48 họ của giai đoạn đầu. Hai tập kiểm thử được gộp vào vì 40 ca vượt qua trên
chúng đã được dùng để thiết kế \ours, nên dùng lại để kiểm định sẽ cho kết quả lạc quan. \Dev{} được chia theo họ thành
phần huấn luyện 126 họ, dùng để chỉnh danh sách từ, ngưỡng và tên trường, và phần kiểm định 42 họ, dùng để chọn cấu hình.
\Test{} gồm 48 họ mới, 8 họ cho mỗi cơ chế, với 560 ca (240 tấn công, 240 \bPlain, 80 \bAuth). Các họ được sinh tự động
từ một đặc tả rồi được hai thành viên của nhóm rà soát câu chữ. Tập được niêm phong bằng SHA-256 trước khi viết dòng mã
đầu tiên của \ours, với độ trùng 4-gram Jaccard so với \dev{} tối đa là 0,05.

\paragraph{Tấn công thích nghi.} Bộ thứ nhất gồm 24 ca thuộc \test, được sinh và niêm phong cùng \test, nhắm vào ba
điểm yếu của \strSand{} với 8 ca cho mỗi điểm yếu. Bộ thứ hai gồm 42 ca do hai thành viên của nhóm soạn sau khi \ours{}
đã đóng băng, với hiểu biết đầy đủ về thiết kế và mô hình hình thức. Mỗi ca thay văn bản tấn công của một họ thuộc phần
kiểm định của \dev{} và nhắm vào các giả định (A2) đến (A5).

\paragraph{Mô hình nền và cấu hình.} Gemma-4-12B, Llama-3.1-8B với ngữ cảnh 16k và Qwen3.5-9B được phục vụ bằng vLLM với
phiên bản cố định, dùng chung prompt ReAct, nhiệt độ 0, tối đa 512 token mỗi bước và tám bước. Gemini 3 Flash của Google
là mô hình đóng nên chỉ được dùng làm cột tham khảo với năm cấu hình. Mọi cấu hình dùng chung tác tử, broker và môi
trường, và khi không phòng thủ, broker chỉ kiểm tra schema. Các phòng thủ khác gồm \cSpot~\cite{hines2024spotlighting},
\cSand, \cPG{} dùng Prompt Guard~2~\cite{meta2025promptguard2} để sàng lọc đầu ra công cụ, \cMel{} và \cProg. Hai phòng
thủ sau được cài đặt lại theo mô tả gốc của MELON~\cite{zhu2025melon} và Progent~\cite{shi2025progent}. Luật tại broker
gồm \phrRules, chặn theo danh sách cụm từ giống chỉ thị, cùng \strRules, \strSand, \ours{} và ba bản bỏ một thành phần.
Tám cấu hình của ma trận chính chạy một lần trên mỗi mô hình nền với 560 ca và 24 ca thích nghi, chỉ chạy lại ca gặp lỗi
hạ tầng. Năm cấu hình \cSpot, \cPG, \cMel, \phrRules{} và \strRules{} đã đóng băng từ trước nhưng được chạy bổ sung sau
khi có kết quả chính, nên so sánh với chúng chỉ mang tính khám phá.

\paragraph{Độ đo.} ASR là tỉ lệ ca tấn công gây hại, tính trên 720 ca tấn công của ba mô hình nền. Tiện ích $U$ là tỉ lệ
ca lành tính hoàn thành theo điều kiện tất định của từng nhiệm vụ mà không gây hại, tính trên 960 ca, và $\Delta U$ là
chênh lệch so với khi không phòng thủ. Lỗi mô hình, như đầu ra sai định dạng hoặc bị cắt, được tính là thất bại về tiện
ích, không bao giờ là phòng thủ thành công. Một họ được coi là \emph{bị hại} nếu có ít nhất một ca tấn công gây hại trên
ít nhất một mô hình nền. Đây là đơn vị suy luận chính, vì các ca trong cùng một họ dùng chung nhiệm vụ và môi trường.
Trên \dev, đơn vị tương ứng là cặp họ và mô hình, gồm 504 cặp. Chi phí gồm tỉ lệ chặn nhầm, tức tỉ lệ ca lành tính có ít
nhất một lần cổng từ chối, tỉ lệ cần xác nhận, số token và thời gian mỗi ca.

\paragraph{Giả thuyết đăng ký trước.} Trước khi niêm phong \test, bốn giả thuyết được đăng ký theo thủ tục kiểm định tuần
tự có cổng. H1 là giả thuyết chính duy nhất: \ours{} có ít họ bị hại hơn \strSand, kiểm bằng phép McNemar chính
xác~\cite{mcnemar1947} một phía trên 48 cặp họ với $\alpha = 0{,}05$. H2 đến H4 chỉ được kiểm định khi giả thuyết không
của H1 bị bác bỏ. H2 cho rằng tiện ích của \ours{} không thấp hơn \strSand{} quá 5 điểm. H3 cho rằng bỏ mỗi thành phần
làm tăng số họ bị hại, có hiệu chỉnh Holm. H4 đặt ngưỡng 10\% cho tỉ lệ cần xác nhận, 15\% cho tỉ lệ chặn nhầm, và mức
tăng so với \strSand{} không quá 20\% về token, 15\% về độ trễ.

\paragraph{AgentDojo.} Bộ Banking với 144 cặp nhiệm vụ người dùng và nhiệm vụ chèn và bộ Slack với 105 cặp của AgentDojo
v1.2.2~\cite{debenedetti2024agentdojo} được chạy với khuôn tấn công \code{important\_instructions}, bằng tiếng Anh và
bằng bản dịch tiếng Việt giữ nguyên IBAN, số tiền, tên và địa chỉ. \StrSand{} và \ours{} chạy với Llama-3.1-8B qua một bộ
chuyển đổi khai báo trước, chỉ rõ công cụ nào là sink và đối số nào chứa đích. Trong bộ chuyển đổi, \pOne{} chỉ yêu cầu
đích xuất hiện nguyên token trong $u$ và không nằm trong phạm vi phủ định. Lời nhắc \cSand{} ở đây là \code{repeat\_user\_prompt} có sẵn của AgentDojo, và tiện ích khi không có tấn công được tính trên 16 nhiệm vụ
Banking và 21 nhiệm vụ Slack.
